\begin{thebibliography}{MCon}
\addcontentsline{toc}{chapter}{Referencias}
\bibitem[Val]{valencia2}
Miguel Ángel Valencia Bucio. \emph{Análisis Matemático 2. Licenciatura en Física y Matemáticas}. Lecturas 1--33, apuntes de curso, sin fecha indicada. Archivos \texttt{AnM2-LFM-lect01.pdf} a \texttt{AnM2-LFM-lect33.pdf}, proporcionados por el autor de esta edición mediante su archivo de trabajo. \href{https://drive.google.com/drive/folders/1HCg3CqIJdUGJjejVX_C304OPpPtwAu44}{Carpeta de las lecturas originales en Google Drive} (acceso sujeto a permisos). Fuente principal del temario y de los ejercicios; la organización, las ampliaciones y las correcciones de esta edición se distinguen de los originales.

\bibitem[Rud]{rudin1987}
Walter Rudin. \emph{Real and Complex Analysis}. Tercera edición. McGraw--Hill, 1987. ISBN 978-0-07-054234-1. \href{https://search.worldcat.org/title/13093736}{Registro bibliográfico de WorldCat}.

\bibitem[Fol]{folland1999}
Gerald B. Folland. \emph{Real Analysis: Modern Techniques and Their Applications}. Segunda edición. John Wiley \& Sons, 1999. ISBN 978-0-471-31716-6. \href{https://www.wiley.com/en-us/Real+Analysis\%3A+Modern+Techniques+and+Their+Applications\%2C+2nd+Edition-p-9780471317166}{Registro editorial de Wiley}.

\bibitem[Bog]{bogachev2007}
Vladimir I. Bogachev. \emph{Measure Theory}. Dos volúmenes. Springer, 2007. \href{https://doi.org/10.1007/978-3-540-34514-5}{DOI: 10.1007/978-3-540-34514-5}.

\bibitem[Hal]{halmos1950}
Paul R. Halmos. \emph{Measure Theory}. Publicación original: D. Van Nostrand, 1950. Reimpresión de Springer, Graduate Texts in Mathematics, vol.~18, 1974. \href{https://doi.org/10.1007/978-1-4684-9440-2}{DOI de la edición electrónica: 10.1007/978-1-4684-9440-2}. La fecha de publicación original se distingue de la incorporación a la colección de Springer.

\bibitem[MCar]{maximenkoCaratheodory}
Egor Maximenko. \emph{Teorema de Carathéodory}. Apunte docente, Escuela Superior de Física y Matemáticas, Instituto Politécnico Nacional, sin fecha indicada, 6 páginas. \href{https://esfm.egormaximenko.com/analysis/Caratheodory_extension_theorem_es.pdf}{PDF en el sitio del autor}. Consultado el 18 de octubre de 2026. Orientación didáctica para separar medida exterior, criterio de medibilidad y extensión de una premedida.

\bibitem[MCon]{maximenkoConvergencias}
Egor Maximenko. \emph{Ejemplos de análisis de varios tipos de convergencia}. Apunte docente, sin fecha indicada, 5 páginas. \href{https://esfm.egormaximenko.com/analysis/convergence_examples_es.pdf}{PDF en el sitio del autor}. Consultado el 18 de octubre de 2026. Contiene ejemplos desarrollados y ejercicios propuestos; la presente edición desarrolla el ejemplo de concentración gaussiana y corrige la variable de límite indicada en su estimación casi uniforme.

\bibitem[MFal]{maximenkoFalsas}
Egor Maximenko. \emph{Algunas demostraciones falsas del teorema de la convergencia monótona (un tema de análisis real)}. Presentación docente, 10 de junio de 2024, 46 páginas de archivo, incluidas superposiciones. \href{https://esfm.egormaximenko.com/analysis/monotone_convergence_theorem_false_proofs_presentation_es.pdf}{PDF en el sitio del autor}. Consultado el 18 de octubre de 2026. Material para discutir cuantificadores y errores de prueba; no se presenta como una demostración completa del teorema.

\bibitem[MLp]{maximenkoLp}
Egor Maximenko. \emph{Los espacios $L^p$ son completos}. Apunte docente, sin fecha indicada, 4 páginas. \href{https://esfm.egormaximenko.com/analysis/Lp_is_complete_es.pdf}{PDF en el sitio del autor}. Consultado el 18 de octubre de 2026. Trata el intervalo $1\le p<\infty$ mediante sucesiones regulares de Cauchy, sumas de diferencias y convergencia monótona.
\end{thebibliography}
